\section*{الفصل \ref{ch:g12:matrix} --- المصفوفات والبيانات}

\begin{solution}{exo:g12:matrix:1}
\[
A + B = \begin{pmatrix} 3 & 2\\ 1 & 2\end{pmatrix}, \quad
AB = \begin{pmatrix} 4 & 2\\ 1 & 1\end{pmatrix}, \quad
BA = \begin{pmatrix} 2 & 4\\ 1 & 3\end{pmatrix}, \quad
A^2 = \begin{pmatrix} 1 & 4\\ 0 & 1\end{pmatrix}.
\]
ولاحظ أن $AB \neq BA$.
\end{solution}

\begin{solution}{exo:g12:matrix:2}
$\det A = 6 - 5 = 1 \neq 0$:
$A^{-1} = \begin{pmatrix} 3 & -5\\ -1 & 2 \end{pmatrix}$.
و $\det B = 12 - 12 = 0$: إذن $B$ ليست \omterm{def:g12:matrix:inverse}{قابلة للقلب}.
\end{solution}

\begin{solution}{exo:g12:matrix:3}
الجملة هي $AX = Y$ حيث $A$ كما في \cref{exo:g12:matrix:2} و
$Y = \begin{pmatrix} 1\\ 2\end{pmatrix}$:
\[
X = A^{-1}Y = \begin{pmatrix} 3 & -5\\ -1 & 2\end{pmatrix}
\begin{pmatrix} 1\\ 2\end{pmatrix}
= \begin{pmatrix} -7\\ 3\end{pmatrix}:
\qquad x = -7,\ y = 3 .
\]
\end{solution}

\begin{solution}{exo:g12:matrix:4}
\emph{1.} $N^2 = \begin{pmatrix} 0&1\\0&0\end{pmatrix}
\begin{pmatrix} 0&1\\0&0\end{pmatrix} = 0$، و $I$ تتبادل مع كل
\omterm{def:g12:matrix:matrix}{مصفوفة}.

\emph{2.} وبما أن الحدين متبادلان، تنطبق مبرهنة ثنائي الحد وتنعدم كل
الحدود التي تحوي $N^2$:
\[
A^k = (2I + N)^k = 2^k I + k\,2^{k-1} N
= \begin{pmatrix} 2^k & k\,2^{k-1}\\ 0 & 2^k \end{pmatrix}.
\]
وللتحقق من أجل $k = 2$:
$A^2 = \begin{pmatrix} 2&1\\0&2\end{pmatrix}^2
= \begin{pmatrix} 4&4\\0&4\end{pmatrix}$، وتعطي الصيغة
$2^2 = 4$ و $2 \times 2 = 4$. \checkmark
\end{solution}

\begin{solution}{exo:g12:matrix:5}
\emph{بالتراجع.} من أجل $n = 1$:
$A^1 = \begin{pmatrix} 0&1\\1&1\end{pmatrix}
= \begin{pmatrix} F_0 & F_1\\ F_1 & F_2\end{pmatrix}$. ولنفترض الصيغة
من أجل $n$؛ عندئذٍ
\[
A^{n+1} = A^n A
= \begin{pmatrix} F_{n-1} & F_n\\ F_n & F_{n+1}\end{pmatrix}
\begin{pmatrix} 0 & 1\\ 1 & 1\end{pmatrix}
= \begin{pmatrix} F_n & F_{n-1} + F_n\\ F_{n+1} & F_n + F_{n+1}\end{pmatrix}
= \begin{pmatrix} F_n & F_{n+1}\\ F_{n+1} & F_{n+2}\end{pmatrix}.
\]
\emph{والمتطابقة.} من أجل المصفوفات $2\times2$، يبيّن \omterm{def:g10:algebra:expand}{النشر} أن
$\det(MN) = \det M \det N$؛ ومنه
$\det(A^n) = (\det A)^n = (-1)^n$، و
$\det A^n = F_{n-1}F_{n+1} - F_n^2$. (وهذه هي \emph{متطابقة كاسيني}.)
\end{solution}

\begin{solution}{exo:g12:matrix:6}
\emph{1.} \omterm{def:g10:coordgeom:system}{بترتيب} الرؤوس $1, 2, 3$:
\[
M = \begin{pmatrix} 0&1&1\\ 0&0&1\\ 1&0&0\end{pmatrix}, \quad
M^2 = \begin{pmatrix} 1&0&1\\ 1&0&0\\ 0&1&1\end{pmatrix}, \quad
M^3 = \begin{pmatrix} 1&1&1\\ 1&0&1\\ 1&0&1 \end{pmatrix}.
\]

\emph{2.} $\bigl(M^3\bigr)_{11} = 1$: أي مسلك مغلق واحد بالضبط طوله $3$
عند الرأس $1$، وهو $1 \to 2 \to 3 \to 1$. (فالمسلك
$1 \to 3 \to 1$ طوله $2$ فقط، و $1 \to 3$ ثم $3\to1$ ثم
$1\to3$ ينتهي عند $3$.)
\end{solution}

\begin{solution}{exo:g12:matrix:7}
\emph{1.} $a_{n+1} = 0.8a_n + 0.3b_n$ و $b_{n+1} = 0.2a_n + 0.7b_n$:
$M = \begin{pmatrix} 0.8 & 0.3\\ 0.2 & 0.7\end{pmatrix}$.

\emph{2.} $MX = X$ يعطي $0.8a + 0.3b = a$، \emph{أي} $0.3b = 0.2a$، إذن
$b = \frac23 a$؛ ومع $a + b = 1$: $a = 0.6$ و $b = 0.4$.

\emph{3.} باستعمال $b_n = 1 - a_n$:
$a_{n+1} = 0.8a_n + 0.3(1 - a_n) = 0.5a_n + 0.3$، إذن
\[
c_{n+1} = a_{n+1} - 0.6 = 0.5a_n + 0.3 - 0.6 = 0.5(a_n - 0.6) = 0.5\,c_n .
\]
ومنه $c_n = 0.5^n c_0 \to 0$: أي $a_n \to 0.6$ و $b_n \to 0.4$، مهما كان
\omterm{def:g11:prob:rv}{التوزيع} الابتدائي.
\end{solution}

\begin{solution}{exo:g12:matrix:8}
\emph{1.} $\det P = -2$، إذن
$P^{-1} = -\frac12\begin{pmatrix} -1 & -1\\ -1 & 1\end{pmatrix}
= \frac12\begin{pmatrix} 1 & 1\\ 1 & -1\end{pmatrix}$. ثم
\[
AP = \begin{pmatrix} 4 & 2\\ 4 & -2 \end{pmatrix}, \qquad
D = P^{-1}AP = \frac12\begin{pmatrix} 1&1\\1&-1\end{pmatrix}
\begin{pmatrix} 4&2\\4&-2\end{pmatrix}
= \begin{pmatrix} 4 & 0\\ 0 & 2\end{pmatrix}.
\]

\emph{2.} من $A = PDP^{-1}$، يعطي تراجع مباشر
$A^n = PD^nP^{-1}$ حيث
$D^n = \begin{pmatrix} 4^n & 0\\ 0 & 2^n\end{pmatrix}$، إذن
\[
A^n = P D^n P^{-1}
= \begin{pmatrix} 4^n & 2^n\\ 4^n & -2^n\end{pmatrix}\cdot
\frac12\begin{pmatrix} 1&1\\1&-1\end{pmatrix}
= \frac12\begin{pmatrix} 4^n + 2^n & 4^n - 2^n\\
4^n - 2^n & 4^n + 2^n\end{pmatrix}.
\]
وتطبيق $A^n$ على $X_0 = \begin{pmatrix} u_0\\v_0\end{pmatrix}$ يعيد
بالضبط صيغ \cref{ex:g12:matrix:coupled}.
\end{solution}

\begin{solution}{pb:g12:matrix:1}
\textbf{1.} $AB = \begin{pmatrix} 2 & 1\\ 4 & 3\end{pmatrix}$
و $BA = \begin{pmatrix} 3 & 4\\ 1 & 2\end{pmatrix}$: فضرب المصفوفات
ليس تبديليًا --- إذ تبادل $B$ الأعمدة عند الضرب من
اليمين، والسطور عند الضرب من اليسار.

\textbf{2.} \omterm{def:g11:vect:det}{المحدد} $6 - 5 = 1$: فالمقلوب
$\begin{pmatrix} 3 & -1\\ -5 & 2\end{pmatrix}$. وبتطبيقه على
$\binom{4}{7}$: $x = 12 - 7 = 5$ و $y = -20 + 14 = -6$.

\textbf{3.} $N^2 = 0$. عندئذٍ $(I + N)^n = I + nN$ بالتراجع:
$(I + nN)(I + N) = I + (n+1)N + nN^2 = I + (n+1)N$.

\textbf{4.} $A = \begin{pmatrix} 0&1&1\\ 1&0&1\\ 1&1&0
\end{pmatrix}$، ومدخلات قطر $A^3$ تساوي $2$: فمن كل
رأس، يوجد بالضبط مسلكان مغلقان طولهما $3$ (وهما المثلث
مسلوكًا في اتجاه عقارب الساعة أو عكسه) --- أي مبرهنة العدّ
في العمل.

\textbf{5.} $D^n = \begin{pmatrix} 2^n & 0\\ 0 & 2^{-n}
\end{pmatrix}$: فاتجاه ينفجر والآخر يموت ---
فالمصائر القطرية متتاليات \omterm{def:g12:seq:arith-geom}{هندسية} مستقلة.

\textbf{6.} $F^2 = \begin{pmatrix} 2 & 1\\ 1 & 1\end{pmatrix}$ و
$F^3 = \begin{pmatrix} 3 & 2\\ 2 & 1\end{pmatrix}$ و
$F^4 = \begin{pmatrix} 5 & 3\\ 3 & 2\end{pmatrix}$: ففيبوناتشي
في كل مكان؛ والتخمين كما صيغ.

\textbf{7.} إذا كان $F^n = \begin{pmatrix} F_{n+1} & F_n\\ F_n &
F_{n-1}\end{pmatrix}$، فإن
\[
F^{n+1} = F^n F =
\begin{pmatrix} F_{n+1} + F_n & F_{n+1}\\
F_n + F_{n-1} & F_n \end{pmatrix}
= \begin{pmatrix} F_{n+2} & F_{n+1}\\ F_{n+1} & F_n
\end{pmatrix} :
\]
وهي الوراثة؛ وأما الحالة الابتدائية $n = 1$ فهي $F$ نفسها، باستعمال
الاصطلاح $F_0 = 0$ (الذي يمدّد العلاقة التراجعية إلى الوراء).

\textbf{8.} $\det F = -1$، إذن
$\det(F^n) = (\det F)^n = (-1)^n$؛ ومباشرةً
$\det(F^n) = F_{n+1}F_{n-1} - F_n^2$: أي كاسيني، في سطر واحد. (وقاعدة
جداء المحددات من أجل المصفوفات $2 \times 2$ نشرٌ لطيف
في خمس دقائق.)

\textbf{9.} أعلى يمين $F^m F^n$:
$F_{m+1}F_n + F_m F_{n-1}$؛ وأعلى يمين $F^{m+n}$:
هو $F_{m+n}$. ومن أجل $m = n = 3$:
$F_4 F_3 + F_3 F_2 = 3 \times 2 + 2 \times 1 = 8 = F_6$.

\textbf{10.} من أجل $k = 1$: بديهي. وإذا كان $F_n \mid F_{kn}$، فإن
صيغة الجمع مع $m = kn$:
$F_{(k+1)n} = F_{kn+1}F_n + F_{kn}F_{n-1}$: فالحدان مضاعفان
للعدد $F_n$. إذن $F_n \mid F_{kn}$ من أجل كل $k$: وللتحقق
$F_3 = 2$ \omterm{def:g12:arith:divides}{يقسم} $F_6 = 8$ و $F_9 = 34$.

\textbf{11.} $F^{100} = F^{64} F^{32} F^4$: أي سبعة تربيعات
($F^2, F^4, \dots, F^{64}$) زائد تركيبين --- تسعة
ضروب بدل تسعة وتسعين. وهي
حيلة جدول المضاعفة عند الكتبة المصريين القدماء، مرفوعةً من
الأعداد إلى المصفوفات: اكتب $100$ بالثنائي، واضرب
المضاعفات التي تحتاج إليها.

\textbf{12.} $0.8 + 0.2 = 1$ و $0.4 + 0.6 = 1$: فالغد
لا بد أن يكون طقسًا \emph{ما} --- فكل عمود \omterm{def:g11:prob:rv}{توزيع}
احتمالي كامل، ومنه فالاحتمالات محفوظة.

\textbf{13.} الغد: $\binom{0.8}{0.2}$. وبعد الغد:
$M\binom{0.8}{0.2} = \binom{0.72}{0.28}$.

\textbf{14.} $Mv = v$ حيث $v = \binom{s}{r}$ و $s + r = 1$:
فإن $0.8s + 0.4r = s$ يعطي $0.4r = 0.2s$، أي $s = 2r$:
$v = \binom{2/3}{1/3}$. فعلى \omterm{def:g10:stats:quartiles}{المدى} الطويل، يومان من كل ثلاثة
مشمسان --- مهما كان شكل اليوم.

\textbf{15.} من $\binom01$: مركّبات الشمس $0.4$ و $0.56$
و $0.624$ و $0.6496$؛ والفجوات إلى $\frac23$: $0.267$ و $0.107$
و $0.043$ و $0.017$ --- فكل خطوة تضرب الفجوة في
$0.4$ بالضبط (وهي القيمة الذاتية الثانية للآلة): أي تقارب
هندسي إلى الحالة المستقرة.

\textbf{16.} الأعمدة (من $A$ و $B$ و $C$):
$P = \begin{pmatrix} 0 & 0 & 1\\ \frac12 & 0 & 0\\
\frac12 & 1 & 0\end{pmatrix}$. والحالة المستقرة: $v_A = v_C$ و
$v_B = \frac{v_A}{2}$ و $v_C = \frac{v_A}{2} + v_B$؛ ومجموعها
$1$: $v = \left(\frac25, \frac15, \frac25\right)$. \omterm{def:g10:coordgeom:system}{والترتيب}:
$A$ و $C$ متعادلتان في الأول، و $B$ أخيرة.

\textbf{17.} تتلقى $C$ \emph{كل} حركة $B$ ونصف
حركة $A$، وهي تصبّ كل شيء عائدًا إلى $A$: فالحالة
المستقرة تقيس أين \emph{يقضي} المتصفح وقته، لا كم
وصلة تشير إليه --- فوصلة واحدة من صفحة رائجة ترجح
عدة وصلات من صفحات مهجورة. وذلك الترجيح التراجعي هو
بالضبط الفكرة المؤسِّسة لغوغل؛ ويعالج التخميد مصائد العناكب
والطرق المسدودة.

\textbf{18.} $A^n = P D^n P^{-1}$ يعطي
$u_n = \frac{4^n + 2^n}{2}$ (و
$v_n = \frac{4^n - 2^n}{2}$). وللتحقق: $u_1 = 3$ و $u_2 = 10$
و $u_3 = 36$؛ ومباشرةً: $(1,0) \to (3,1) \to (10,6) \to
(36, 28)$: فهما متطابقان.

\textbf{19.} تغيّر القطرنة \omterm{def:g10:coordgeom:system}{الإحداثيات} إلى \omterm{def:g10:coordgeom:system}{إحداثيات} تتفكك فيها
الجملة المقترنة إلى متتاليات \omterm{def:g12:seq:arith-geom}{هندسية} مستقلة ---
فكل قيمة ذاتية تجري سباقها. والحالة المستقرة لسلسلة ماركوف
هي \omterm{def:g11:vect:vector}{المتجهة} الذاتية للقيمة الذاتية $1$، ومعدل
تقارب السؤال 15 هو القيمة الذاتية التالية: فآلة
الطقس كانت حكاية ذاتية منذ البداية.

\textbf{20.} مسك الدفاتر: فالجملة \omterm{def:g10:algebra:equation}{معادلة} مصفوفية واحدة،
تُحل بمقلوب واحد. والتوافيق: فقوى \omterm{def:g12:matrix:graph}{مصفوفة الجوار}
تعدّ المسالك والوصلات والاتصالات. والتطور: فقوى
الآلة تحمل الحالات إلى مصائرها، و
الاتجاهات الذاتية (اتجاه فيبوناتشي الذهبي، والحالة المستقرة
للطقس، \omterm{def:g11:vect:vector}{ومتجهة} \omterm{def:g10:coordgeom:system}{ترتيب} الشبكة) هي المصائر.
والجبر الخطي، في الكتب الجامعية، هو علم هذا بالضبط.
\end{solution}
